\documentclass[10pt]{amsart}
\usepackage[a4paper,margin=23mm]{geometry}
\usepackage{amsmath,amssymb,amsthm,mathtools,hyperref,graphicx,xspace,letltxmacro,xpatch}
\usepackage[english]{babel}
\newtheorem{theo}{Theorem}[section]
\newtheorem{lemm}[theo]{Lemma}\newtheorem{prop}[theo]{Proposition}\newtheorem{coro}[theo]{Corollary}\newtheorem{conj}[theo]{Conjecture}
\theoremstyle{definition}\newtheorem{defi}[theo]{Definition}\newtheorem{rema}[theo]{Remark}\newtheorem{exam}[theo]{Example}\newtheorem{exem}[theo]{Example}\newtheorem{nota}[theo]{Notation}
\newcommand{\ie}{i.e.}\newcommand{\eg}{e.g.}\newcommand{\resp}{resp.}\newcommand{\skpt}{\par\smallskip\noindent}
\emergencystretch=3em\pagestyle{plain}
%~Mouliné par MaN_auto v.0.23.0 2020-05-25 16:49:39
%
%\makeatletter
%%\def\@altabstractlanguage{}
%\makeatother


\usepackage{tikz}
\usepackage{tkz-tab}
\usepackage{array}
%\usepackage[normalem]{ulem}
%\usepackage{multirow}
%\usepackage{amscd}
\usepackage{mathtools}

\DeclarePairedDelimiter{\vect}{\langle}{\rangle}

\let\oldtilde\tilde
\renewcommand*{\tilde}[1]{\mathchoice{\widetilde{#1}}{\widetilde{#1}}{\oldtilde{#1}}{\oldtilde{#1}}}
\let\oldhat\hat
\renewcommand*{\hat}[1]{\mathchoice{\widehat{#1}}{\widehat{#1}}{\oldhat{#1}}{\oldhat{#1}}}

\makeatletter
\let\save@mathaccent\mathaccent
\newcommand*\if@single[3]{%
  \setbox0\hbox{${\mathaccent"0362{#1}}^H$}%
  \setbox2\hbox{${\mathaccent"0362{\kern0pt#1}}^H$}%
  \ifdim\ht0=\ht2 #3\else #2\fi
  }
%The bar will be moved to the right by a half of \macc@kerna, which is computed by amsmath:
\newcommand*\rel@kern[1]{\kern#1\dimexpr\macc@kerna}
%If there's a superscript following the bar, then no negative kern may follow the bar;
%an additional {} makes sure that the superscript is high enough in this case:
\newcommand*\widebar[1]{\@ifnextchar^{{\wide@bar{#1}{0}}}{\wide@bar{#1}{1}}}
%Use a separate algorithm for single symbols:
\newcommand*\wide@bar[2]{\if@single{#1}{\wide@bar@{#1}{#2}{1}}{\wide@bar@{#1}{#2}{2}}}
\newcommand*\wide@bar@[3]{%
  \begingroup
  \def\mathaccent##1##2{%
%Enable nesting of accents:
    \let\mathaccent\save@mathaccent
%If there's more than a single symbol, use the first character instead (see below):
    \if#32 \let\macc@nucleus\first@char \fi
%Determine the italic correction:
    \setbox\z@\hbox{$\macc@style{\macc@nucleus}_{}$}%
    \setbox\tw@\hbox{$\macc@style{\macc@nucleus}{}_{}$}%
    \dimen@\wd\tw@
    \advance\dimen@-\wd\z@
%Now \dimen@ is the italic correction of the symbol.
    \divide\dimen@ 3
    \@tempdima\wd\tw@
    \advance\@tempdima-\scriptspace
%Now \@tempdima is the width of the symbol.
    \divide\@tempdima 10
    \advance\dimen@-\@tempdima
%Now \dimen@ = (italic correction / 3) - (Breite / 10)
    \ifdim\dimen@>\z@ \dimen@0pt\fi
%The bar will be shortened in the case \dimen@<0 !
    \rel@kern{0.6}\kern-\dimen@
    \if#31
      \overline{\rel@kern{-0.6}\kern\dimen@\macc@nucleus\rel@kern{0.4}\kern\dimen@}%
      \advance\dimen@0.4\dimexpr\macc@kerna
%Place the combined final kern (-\dimen@) if it is >0 or if a superscript follows:
      \let\final@kern#2%
      \ifdim\dimen@<\z@ \let\final@kern1\fi
      \if\final@kern1 \kern-\dimen@\fi
    \else
      \overline{\rel@kern{-0.6}\kern\dimen@#1}%
    \fi
  }%
  \macc@depth\@ne
  \let\math@bgroup\@empty \let\math@egroup\macc@set@skewchar
  \mathsurround\z@ \frozen@everymath{\mathgroup\macc@group\relax}%
  \macc@set@skewchar\relax
  \let\mathaccentV\macc@nested@a
%The following initialises \macc@kerna and calls \mathaccent:
  \if#31
    \macc@nested@a\relax111{#1}%
  \else
%If the argument consists of more than one symbol, and if the first token is
%a letter, use that letter for the computations:
    \def\gobble@till@marker##1\endmarker{}%
    \futurelet\first@char\gobble@till@marker#1\endmarker
    \ifcat\noexpand\first@char A\else
      \def\first@char{}%
    \fi
    \macc@nested@a\relax111{\first@char}%
  \fi
  \endgroup
}
\makeatother
\let\oldbar\bar
\renewcommand*{\bar}[1]{\mathchoice{\widebar{#1}}{\widebar{#1}}{\oldbar{#1}}{\oldbar{#1}}}

%%%%%%%%%%%%%%


\newcounter{a}
\setcounter{a}{0}
\newcommand{\Cond}{Cond.}
\newenvironment{cond}{\refstepcounter{a}{\upshape\emph{Cond.~\thea}}\pointir}{}
\newenvironment{condd}{{\upshape\emph{Cond~\thea*}\pointir}}{}
\newenvironment{conddd}[1]{{\upshape\emph{Cond~\ref{#1}}\pointir}}{}

%\theoremstyle{remark}
%\newtheorem{cond}{Cond.}
%\newtheorem{condd}{Cond.}
%\renewcommand*{\thecondd}{\thecond*}


\newcommand{\tT}{\tilde T}

\newcommand{\encadre}[1]{\begin{tabular}{|c|} \hline #1 \\ \hline\end{tabular}}
% Ça marchait pas trop mal, mais je trouve que ma version est un peu mieux
\newcommand{\myencadre}[2][5.1cm]{\vspace{0.15cm}\fbox{\parbox{#1}{#2}}\vspace{0.15cm}}
 
\newcommand{\app}[5]{\begin{array}{rccl} #1:&#2&\longrightarrow &#3\\ &#4&\longmapsto &#5\end{array}}
%\newcommand{\rv}[1]{\mathbf{#1}}

%% Commande ci dessous désactivée en raison des problèmes d'alignement causées dans les formules :
%%remplacement en manuel
%%\newcommand{\prob}[1]{\mathbb{P}\left(#1\right)}

%%COMMANDES OK
\newcommand{\pw}[1]{\mathbb{P}_W\left(#1\right)}
\newcommand{\eqd}{\stackrel{(d)}{=}}
\newcommand{\ind}[1]{1_{#1}}
\newcommand{\diff}[1]{\mathrm{d} #1}
%%\DeclareMathOperator{\diff}{d}
\newcommand{\borel}[1]{\mathcal{B}\left(#1 \right)}
\newcommand{\triang}[2]{\mathcal{T}_{#1} \left(#2 \right)}
\newcommand{\triangE}[1]{\mathcal{T}_{#1}}



%\newcommand{\id}{\mathrm{id}}%%mathop
\DeclareMathOperator{\id}{id}

%%math caract%%%%%%%%%%%%%%%%
\newcommand{\NN}{\mathbb{N}}
\newcommand{\ZZ}{\mathbb{Z}}
\newcommand{\Z}{\mathbb{Z}}
\newcommand{\Ze}{\mathbb{Z}_{\mathrm{e}^2}}
\newcommand{\RR}{\mathbb{R}}
\newcommand{\Pbb}{\mathbb{P}}

\newcommand{\ber}{\mathcal{B}}
\renewcommand{\L}{\mathcal{L}}

%%Mathop%%%%%%%%%%%%%%%%%%%%%%%%%%%%

%%\newcommand{\vu}{{\bf \textup{u}}}
\DeclareMathOperator{\vu}{u}
%%\newcommand{\vv}{{\bf \textup{v}}}
\DeclareMathOperator{\vv}{v}

\newcommand{\Fh}{M_{-\vv}}
\newcommand{\Bh}{M_{\vu}}

%%\newcommand{\sym}{\text{sym}}
\DeclareMathOperator{\sym}{sym}


%\newcommand{\M}{\mathcal{M}}

%\newcommand{\ppm}{\Pi_p}

\newcommand{\haut}[1]{$\begin{array}{c}(#1)\\ \;\end{array}$}
\newcommand{\espace}{\rule{0pt}{6ex}}
\newcommand{\test}{\vspace{0.15cm}\rule{0pt}{5ex}\vspace{0.15cm}}

%\renewcommand{\thefootnote}{\alph{footnote}}

\newcommand{\isdef}{\triangleq}


%%%%%%%%%%%%%%%%%%%%%%%%%%%%%%%%%%
\definecolor{grey}{rgb}{0.5,0.5,0.5}

\usetikzlibrary{arrows}
\usetikzlibrary{decorations.markings}
\tikzset{->-/.style={decoration={ markings, mark=at position.7 with {\arrow[scale=2]{>}}},postaction={decorate}}}
\tikzset{>>/.style={decoration={ markings, mark=at position.7 with {\arrow[scale=5]{>}}},postaction={decorate}}}
\tikzset{>>>/.style={decoration={ markings, mark=at position 1 with {\arrow[scale=5]{>}}},postaction={decorate}}}

%%%%%%%%%%%%%%%%%%%%%%%%%%%%%%%%%%%%%%%%%%

\graphicspath{{./figures/}}

\newcommand*{\mk}{\mkern -1mu}
\newcommand*{\Mk}{\mkern -2mu}
\newcommand*{\mK}{\mkern 1mu}
\newcommand*{\MK}{\mkern 2mu}

\hypersetup{urlcolor=purple, linkcolor=blue, citecolor=red}

\newcommand*{\romanenumi}{\renewcommand*{\theenumi}{\roman{enumi}}}
\newcommand*{\Romanenumi}{\renewcommand*{\theenumi}{\Roman{enumi}}}
\newcommand*{\alphenumi}{\renewcommand*{\theenumi}{\alph{enumi}}}
\newcommand*{\Alphenumi}{\renewcommand*{\theenumi}{\Alph{enumi}}}
\let\oldforall\forall
\renewcommand*{\forall}{\mathrel{\oldforall}}
%%%%%%%%%%%%%%%%%%%%%%%%%%%%%%%%%%%%%%%%%%%%%%

\providecommand{\pointir}{.\ }

\newtheorem{definition}[theo]{Definition}
\newtheorem{lemma}[theo]{Lemma}
\newtheorem{proposition}[theo]{Proposition}
\newtheorem{corollary}[theo]{Corollary}
\newtheorem{theorem}[theo]{Theorem}
\newtheorem{remark}[theo]{Remark}
\newenvironment{proofc}[1][\proofname]{\begin{proof}[{#1}]}{\end{proof}}
